% ---- DSA27 handout: XeLaTeX preamble -------------------------------------

% Arabic: XeTeX native right-to-left, without polyglossia/bidi.
% polyglossia pulls in bidi, which patches tabular and breaks pandoc's
% longtable. XeTeX already provides \beginR/\endR, which is all that short
% Arabic runs inside English text need. tools/arabic.lua wraps them.
%
% The Arabic face is passed in by tools/build.py (Segoe UI on Windows, Noto
% Sans Arabic elsewhere). It carries regular and bold only, so the italic slots
% are mapped back onto them -- otherwise Arabic inside emphasis or a blockquote
% silently disappears.
\TeXXeTstate=1
\newfontfamily\arabicfont[
  Script=Arabic, Scale=1.12,
  BoldFont={$arabicboldfont$},
  ItalicFont={$arabicfont$},
  BoldItalicFont={$arabicboldfont$}]{$arabicfont$}

% Symbols the main font may lack (Noto Sans has no arrows, minus, <=, box
% drawing...). Each is set from $symbolfont$, which tools/build.py picks as a
% font known to carry them all, so prose and code render them identically.
\usepackage{newunicodechar}
\newfontfamily\dsasymbolfont{$symbolfont$}
\newunicodechar{→}{{\dsasymbolfont\symbol{"2192}}}
\newunicodechar{−}{{\dsasymbolfont\symbol{"2212}}}
\newunicodechar{─}{{\dsasymbolfont\symbol{"2500}}}
\newunicodechar{≥}{{\dsasymbolfont\symbol{"2265}}}
\newunicodechar{├}{{\dsasymbolfont\symbol{"251C}}}
\newunicodechar{│}{{\dsasymbolfont\symbol{"2502}}}
\newunicodechar{≤}{{\dsasymbolfont\symbol{"2264}}}
\newunicodechar{└}{{\dsasymbolfont\symbol{"2514}}}
\newunicodechar{√}{{\dsasymbolfont\symbol{"221A}}}
\newunicodechar{←}{{\dsasymbolfont\symbol{"2190}}}
\newunicodechar{≈}{{\dsasymbolfont\symbol{"2248}}}
\newunicodechar{┐}{{\dsasymbolfont\symbol{"2510}}}
\newunicodechar{►}{{\dsasymbolfont\symbol{"25BA}}}
\newunicodechar{┘}{{\dsasymbolfont\symbol{"2518}}}
\newunicodechar{∞}{{\dsasymbolfont\symbol{"221E}}}
\newunicodechar{⌊}{{\dsasymbolfont\symbol{"230A}}}
\newunicodechar{⌋}{{\dsasymbolfont\symbol{"230B}}}
\newunicodechar{⇒}{{\dsasymbolfont\symbol{"21D2}}}
\newunicodechar{↔}{{\dsasymbolfont\symbol{"2194}}}
\newunicodechar{⌈}{{\dsasymbolfont\symbol{"2308}}}
\newunicodechar{⌉}{{\dsasymbolfont\symbol{"2309}}}
\newcommand{\textarabic}[1]{{\arabicfont\beginR #1\endR}}

\usepackage{xcolor}
\definecolor{dsaslate}{HTML}{233A3E}
\definecolor{dsaamber}{HTML}{C8860D}
\definecolor{dsarule}{HTML}{D8DCDA}
\definecolor{dsamuted}{HTML}{5B6B6E}

\usepackage{titlesec}
\titleformat*{\section}{\Large\bfseries\color{dsaslate}}
\titleformat*{\subsection}{\large\bfseries\color{dsaslate}}
\titleformat*{\subsubsection}{\normalsize\bfseries\color{dsaslate}}

\usepackage{booktabs}
\usepackage{array}

% Quote blocks get an amber left rule instead of plain indentation.
\usepackage{mdframed}
\renewenvironment{quote}
  {\begin{mdframed}[leftline=true,topline=false,bottomline=false,rightline=false,
     linewidth=2pt,linecolor=dsaamber,backgroundcolor=white,
     innerleftmargin=10pt,innertopmargin=6pt,innerbottommargin=6pt,
     skipabove=8pt,skipbelow=8pt]\itshape}
  {\end{mdframed}}

\usepackage{fancyhdr}
\pagestyle{fancy}
\fancyhf{}
\renewcommand{\headrulewidth}{0.4pt}
\fancyhead[L]{\small\color{dsamuted}DSA27 \textperiodcentered{} Data Structures and Algorithms}
\fancyfoot[C]{\small\color{dsamuted}\thepage}
